% ======================================================================
% Chapter 6 -- Part B: build, deploy, evaluate, business case, end matter
% ======================================================================

\section{Build: How Do We Make It?}
\label{sec:inv-build}

\subsection{The workflow}

The risk-routed design is a graph workflow in ADK~2, built the way the sample builds its expense flow: routing decisions are code functions that return a named route, model calls are nodes, and pauses for people are nodes that yield a request for input. Listing~\ref{lst:inv-workflow} shows the graph and the routing function.

\begin{lstlisting}[language=Python, caption={The risk-routed design's workflow (abridged from \texttt{ch06/invoice\_agent/agent.py}).}, label={lst:inv-workflow}]
from google.adk import Agent, Context, Event, Workflow
from google.adk.events import RequestInput
from .controls import run_controls          # match, duplicates, bank check, score
from .config import config

def check_invoice(node_input: dict, ctx: Context) -> Event:
    """Run the deterministic controls and choose the route."""
    ctx.state["invoice"] = node_input
    checks = run_controls(node_input)
    ctx.state["checks"] = checks
    if checks["bank_details_changed"]:
        return Event(route="TREASURY", output=node_input)
    if checks["clean_match"] and checks["risk"] < config.review_threshold:
        return Event(route="AUTO", output=node_input)
    return Event(route="REVIEW", output=node_input)

def clerk_review(node_input, ctx: Context):
    """Pause until a clerk resolves the invoice in the review interface."""
    yield RequestInput(
        message="Invoice needs review.",
        payload={"invoice": ctx.state["invoice"],
                 "checks": ctx.state["checks"],
                 "evidence": node_input},          # from the explanation agent
    )

root_agent = Workflow(
    name="invoice_to_payment",
    edges=[
        ("START", parse_event, extraction_agent, validate_invoice),
        (validate_invoice, {"VALID": check_invoice, "INVALID": clerk_review}),
        (check_invoice, {"AUTO": approval_by_matrix,
                         "REVIEW": explanation_agent,
                         "TREASURY": treasury_review}),
        (explanation_agent, clerk_review, record_clerk_decision),
        (record_clerk_decision, {"RESOLVED": approval_by_matrix,
                                 "REJECTED": notify_supplier}),
        (approval_by_matrix, {"NO_APPROVAL_NEEDED": schedule_payment,
                              "NEEDS_APPROVAL": budget_holder_review}),
        (budget_holder_review, record_approval, schedule_payment),
    ],
)
\end{lstlisting}

The structure makes three properties easy to verify by reading: every route out of \texttt{check\_invoice} is named; the only path to \texttt{schedule\_payment} passes either the authority matrix or a recorded approval; and the model nodes (\texttt{extraction\_agent}, \texttt{explanation\_agent}) appear nowhere on the routing decisions.

\subsection{Designing the checkpoint}

The technical focus of this case is the human checkpoint, and a checkpoint is more than a pause. Five questions define it.

\begin{thinkabout}
Imagine you are the clerk. An invoice arrives flagged as a possible duplicate. What would you need to see on one screen to decide in under a minute, and what would make you approve it without really looking?
\end{thinkabout}

\paragraph{What does the reviewer see?} The evidence pack: the invoice image with the extracted fields beside it; the purchase-order and receipt lines that disagree, with the differences highlighted; the result of each control (duplicate candidates with links, bank-detail history); and the supplier's last ten invoices. The explanation agent writes a short summary of \emph{what failed}, citing the lines; it does not recommend approval or rejection. A recommendation from the system invites the reviewer to agree with it, which is the opposite of what a reviewer is for.

\paragraph{What may the reviewer decide?} A closed set of decisions, each with required fields: \emph{approve} (with a reason code if any control failed), \emph{reject} (with a reason sent to the supplier), \emph{request information} (from the buyer, the warehouse or the supplier, which restarts the deadline), or \emph{correct} (edit an extracted field, which re-runs the controls). There is no free-text decision.

\paragraph{Who may decide?} The authority matrix---who may approve what amount for which cost centre---is data owned by finance and enforced in code. The requester of a purchase may never approve its invoice. The decision node reads the reviewer's identity from the verified session of the review interface, never from the decision payload.

\paragraph{What if nobody decides?} Each pause has a deadline: one business day for clerks, two for budget holders, or less if the discount date is nearer. A scheduled job finds paused workflows past their deadline and escalates to the reviewer's manager, then to the controller. Nothing is approved by timeout.

\paragraph{Is the reviewer really reviewing?} Reviewers who see mostly clean items learn to approve quickly; this is called \keyterm{rubber-stamping}\index{rubber-stamping|textbf}, and it silently converts a control into a delay. Three measures detect it. \textbf{Canaries}: about 1\% of review items are seeded copies of known problems (a duplicate, an inflated price), and each reviewer's catch rate is tracked. \emph{Time on task}: decisions made in a few seconds on items with failed controls are flagged. \textbf{Agreement}: a sample of each reviewer's decisions is re-reviewed by a senior clerk.

\subsection{Failing closed}

Listing~\ref{lst:strict-parse} replaces the sample's decision parsing. It accepts only a structured decision from the review form, checks it against a closed set, and enforces authority and segregation of duties. Anything else is \emph{invalid}, which keeps the workflow paused and alerts the reviewer, rather than being read as a decision.

\lstinputlisting[language=Python, caption={Fail-closed decision parsing (\texttt{ch06/decision.py}).}, label={lst:strict-parse}, firstline=11, lastline=23]{code/decision.py}

Running the sample's rule and this one side by side on four inputs:

\begin{lstlisting}[language={}, numbers=none, frame=none, xleftmargin=1.5em]
'approve'          -> sample: approve
'Do not approve'   -> sample: approve
'I disapprove'     -> sample: approve
'rejected'         -> sample: reject
('invalid', 'Decision must come from the review form, not free text.')
('invalid', "Amount exceeds this reviewer's approval limit.")
('approve', '')
\end{lstlisting}

\begin{principle}[Fail closed where money moves]
Any step that releases money, grants access or changes a record of authority must treat every unexpected input---malformed, ambiguous, missing or from an unauthorized person---as \emph{no}.
\end{principle}

% ----------------------------------------------------------------------
\section{Deploy: How Do We Run It Safely?}
\label{sec:inv-deploy}

The sample's Terraform provides most of what the risk-routed design needs: a Cloud Run service for the agent, a second for the review interface behind Identity-Aware Proxy, a Pub/Sub topic with a push subscription that triggers the workflow, and log-based metrics with alert policies. The risk-routed design adds five things:
\begin{enumerate}[itemsep=2pt]
  \item \textbf{Durable pauses.} A workflow paused for two days must survive restarts and new deployments. Sessions are stored in a managed database rather than in memory, and the agent service is deployed so that a new version can resume sessions created by the old one (Chapter~\ref{ch:capstone} returns to versioning).
  \item \textbf{Deadlines.} A Cloud Scheduler job runs every hour, finds paused workflows past their deadline, and escalates.
  \item \textbf{Separation of payment.} The payment server belongs to treasury and runs under treasury's service account. The workflow's service account can call \texttt{schedule\_payment} but cannot read or change supplier bank details; no component of the solution can.
  \item \textbf{An audit trail.} Every routing decision, control result, reviewer decision and approval is written as a structured log entry and exported to a write-once storage bucket that auditors can query. Chapter~\ref{ch:case-kyc} develops audit trails in depth.
  \item \textbf{Shadow mode.} A configuration switch makes \texttt{schedule\_payment} record what it \emph{would} have done instead of doing it. The pilot begins in this mode (Section~\ref{sec:inv-business}).
\end{enumerate}

% ----------------------------------------------------------------------
\section{Evaluate: How Do We Know It Works?}
\label{sec:inv-evaluate}

\begingroup
\small
\begin{xltabular}{\textwidth}{@{}>{\raggedright\arraybackslash}X>{\raggedright\arraybackslash}p{3.0cm}r@{}}
\caption{Evaluation thresholds for the risk-routed design.}\label{tab:inv-thresholds}\\
\toprule
\rowcolor{coHead} Metric & Step & Threshold \\
\midrule
\endfirsthead
\toprule
\rowcolor{coHead} Metric & Step & Threshold \\
\midrule
\endhead
Invoice total and supplier extracted correctly (PDFs) & Extraction, validation & $\ge$ 99.5\% \\
Invalid extraction reaching the controls & Validation & 0 \\
Seeded duplicates detected & Controls & 100\% \\
Share of problem \textbf{value} routed to review, at review budget & Risk score & $\ge$ 90\% \\
Bank-detail changes routed to treasury & Router & 100\% \\
Evidence-pack claims supported by the cited lines & Explanation agent & $\ge$ 98\% \\
Approvals outside authority or by the requester & Decision node & 0 \\
Canary catch rate, per reviewer & Clerks & $\ge$ 95\% \\
\bottomrule
\end{xltabular}
\endgroup

Two features of the evaluation are specific to human-in-the-loop systems. First, \emph{the people are part of what is evaluated}: canary catch rates and re-review agreement measure the checkpoint, not just the software. Second, \emph{controls are tested the way auditors test them}: for \risk{Wrong approver}, the test is an attempt by an unauthorized user, and by a requester, to approve through every path, including by editing the payload sent from the review interface.

\begin{keyidea}
A control that has not been tried is an assertion, not a control.
\end{keyidea}

The risk score is evaluated on the share of problem \textbf{value} that reaches review, not on accuracy. A score that is 99\% accurate because 96\% of invoices are clean is useless; what matters is how much of the money at risk it puts in front of a person within the review budget.

% ----------------------------------------------------------------------
\section{Business Case: Is It Worth It?}
\label{sec:inv-business}

\paragraph{Value.} At about \$2.90 per invoice against \$9.40 today, the risk-routed design frees about \$62{,}000 a month of clerk time at current volume. The controller's intended use for that capacity is specific: absorb the acquisition's 40\% increase in volume without hiring the four clerks it would otherwise need (about \$260{,}000 a year). Faster approval raises captured early-payment discounts from 35\% to a target of 85\%: on about \$4.6 million a month of discount-eligible spend at 2\%, that is about \$46{,}000 a month, or \$557{,}000 a year, and it rests on a measured input---which invoices currently miss their discount date and why. Losses stay close to what reviewing every invoice would achieve (Table~\ref{tab:review-policies}), with bank-detail fraud addressed by a control that does not depend on a model at all. Build cost is estimated at \$350{,}000.

\begin{quickcheck}
Which of the value lines above rests on a measured input, and which on an assumption?
\answer{The discount line rests on a measured input: which invoices miss their discount date today, and why. The capacity line depends on the acquisition going ahead and on the controller's decision not to hire.}
\end{quickcheck}

\paragraph{Pilot: shadow, then partial, then full.} For a workflow that moves money, the safest pilot runs in three stages.
\begin{enumerate}[itemsep=1pt]
  \item \emph{Shadow (four weeks).} The risk-routed design processes every invoice in shadow mode while clerks work as today. Every difference between the risk-routed design's route and the clerks' outcome is reviewed. This measures extraction, controls and the risk score on real invoices with no exposure.
  \item \emph{Partial (four weeks).} The touchless path is enabled for invoices from the 200 suppliers with the cleanest history, under \$5{,}000. All other invoices behave as in the perfect-match design: every exception to a clerk.
  \item \emph{Full.} The review threshold moves to the value computed from the shadow data.
\end{enumerate}
\emph{Go} from each stage to the next requires zero duplicate payments, zero unauthorized approvals, at least 99.5\% extraction accuracy on totals, and canary catch rates of at least 95\%. The controller owns each decision; internal audit signs off before the touchless path is enabled.

\begin{tcolorbox}[colback=white, colframe=black!60, boxrule=0.6pt, arc=1pt, left=10pt, right=10pt, top=8pt, bottom=8pt, fontupper=\small]
\textbf{To:} Controller \hfill \textbf{From:} AP Design Team\\
\textbf{Re:} Invoice processing that captures discounts and absorbs the acquisition --- request to pilot

\smallskip\hrule\smallskip

\textbf{Recommendation.} Approve a three-stage pilot starting in shadow mode; build cost about \$350{,}000.

\textbf{The problem.} Invoices take 11 days to approve and cost \$9.40 each to process. We capture 35\% of early-payment discounts. Duplicates, overcharges and two bank-detail frauds last year cost us money. The acquisition adds 40\% to volume.

\textbf{What changes.} Invoices are read and checked automatically against orders and receipts. Clean, low-risk invoices go straight to approval under the existing authority matrix. Clerks see only invoices whose risk justifies their time, with the evidence laid out. Any change to a supplier's bank details always goes to treasury for a call-back check.

\textbf{Value.} About \$557{,}000 a year in additional discounts; the acquisition's volume handled without four new hires (about \$260{,}000 a year); losses kept near the level of reviewing every invoice.

\textbf{Controls.} Approval authority and segregation of duties are enforced in the system. Nothing is approved by timeout. Every decision is recorded for audit. Internal audit signs off before any invoice is paid without a person.

\textbf{What we need from you.} Treasury's ownership of the payment step; internal audit's participation in the shadow stage; agreement to the review threshold being set from pilot data rather than fixed in advance.
\end{tcolorbox}

% ----------------------------------------------------------------------
\section{Changed Requirement: What If the World Changes?}
\label{sec:inv-changed}

\begin{enumerate}[itemsep=3pt]
  \item \textbf{E-invoicing mandate.} Two of Meridian's markets will require structured electronic invoices from suppliers within two years. Which steps become unnecessary for those invoices? Does the case for the risk-routed design over the perfect-match design get weaker?
  \item \textbf{Different systems after the acquisition.} The acquired company runs a different ERP. Should the risk-routed design call both ERPs through one MCP server\index{Model Context Protocol (MCP)!server} or two? Who owns each?
  \item \textbf{Dual approval.} The external auditor requires two approvers for any invoice above \$50{,}000. Change the design table and the decision step. What happens to \kpi{Days to approval} for large invoices, and how would you protect it?
\end{enumerate}

% ----------------------------------------------------------------------
\section*{Chapter Summary}
\addcontentsline{toc}{section}{Chapter Summary}

\begin{itemize}[itemsep=2pt]
  \item The \texttt{ambient-expense-agent} sample shows the right structure---routing in code, models only for judgment, a pause for people---and leaves gaps that matter when money moves: routing by amount, structured inputs assumed, no matching or fraud controls, a decision parser that can fail open, no authority checks, no deadlines.
  \item Review should be placed by risk, not by amount. In Meridian's simulation, risk-scored review checked under 6\% of invoices and lost about as little as reviewing all of them, at a twentieth of the labor.
  \item A checkpoint is defined by five answers: what the reviewer sees, what they may decide, who may decide, what happens when nobody does, and how rubber-stamping is detected.
  \item Fail closed wherever money moves. Never approve by timeout. Take identity from the verified session, never from the payload.
  \item Documents from outside the company sit inside a trust boundary; values that matter, such as bank details, never come from them.
  \item Pilot money-moving workflows in shadow mode first, then on the safest segment, then fully.
\end{itemize}

\section*{Review Questions}
\addcontentsline{toc}{section}{Review Questions}

\begin{enumerate}[itemsep=2pt]
  \item List three strengths of the sample's design and explain why each one carries over to the risk-routed design.
  \item Why does ``Do not approve'' approve in the sample's free-text path? Rewrite the rule so that it fails closed.
  \item Explain why amount thresholds are a poor review rule for invoices, using Table~\ref{tab:review-policies}.
  \item Why does the evidence pack describe what failed but not recommend a decision?
  \item What are canaries, and what do they measure that ordinary accuracy metrics cannot?
  \item Why are bank details never taken from an invoice?
  \item Why is there no ``approve by timeout,'' even when a discount deadline is near?
  \item What does shadow mode measure, and what can it not measure?
\end{enumerate}

\section*{Exercises}
\addcontentsline{toc}{section}{Exercises}

\begin{enumerate}[itemsep=2pt]
  \item \textbf{(Design)} Write the perfect-match design as a design table and run the eight checks of Chapter~\ref{ch:toolkit} on all three designs.
  \item \textbf{(Design)} Write the component card for the explanation agent, including what it must never say.
  \item \textbf{(Design)} Complete one changed requirement from Section~\ref{sec:inv-changed}.
  \item \textbf{(Build)} In \texttt{review\_policy.py}, make review cost rise with invoice complexity and reviewer catch rate fall as daily review volume rises. Recompute the best threshold.
  \item \textbf{(Build)} Implement the deadline sweep: a function that, given paused workflows with deadlines, returns the escalations due, and its tests.
\end{enumerate}

\section*{Lab 6: A Checkpoint That Holds}
\addcontentsline{toc}{section}{Lab 6: A Checkpoint That Holds}

\paragraph{Goal.} Build the risk-routed design from the sample, run it in shadow mode on a synthetic month of invoices, set the review threshold from the results, and show that the checkpoint resists unauthorized and malformed decisions.

\paragraph{Steps.}
\begin{enumerate}[itemsep=2pt]
  \item In \texttt{ch06/}, deploy the pinned sample with its Terraform and send it three test expenses, including one resumed with the text ``do not approve.'' Record the outcome.
  \item Complete \texttt{controls.py} (match, duplicates, bank-detail check) and its tests against the synthetic ERP.
  \item Assemble the workflow of Listing~\ref{lst:inv-workflow} and replace the decision node with Listing~\ref{lst:strict-parse}.
  \item Publish the 9{,}500 synthetic invoices to the Pub/Sub topic in shadow mode. Compute Table~\ref{tab:review-policies} from the workflow's own logs and choose the threshold.
  \item Attempt approval as an unauthorized user, as the requester, and by editing the decision payload. All three must fail.
  \item Run the cleanup script.
\end{enumerate}

\paragraph{Deliverable.} The review-policy table from your run, the chosen threshold with its justification, the control-test results, and a one-paragraph answer to: \emph{which of the sample's gaps would an external auditor find first?}
